\chapter*{Résumé}
\addcontentsline{toc}{chapter}{Résumé}

\selectlanguage{french}
IntelliConnect est une plateforme de gestion des partenariats et des projets d'entreprise conçue et développée pour Capgemini Engineering Tunisie. La plateforme offre à l'entreprise un espace de travail unique et sécurisé pour piloter son réseau de partenaires et son portefeuille de projets, en remplaçant des fichiers dispersés et des outils déconnectés par un seul système gouverné. Elle rassemble les données des partenaires et des projets, suit l'activité et les engagements dans le temps, et aide les équipes à planifier et à suivre le travail qui maintient chaque relation active. Le contrôle d'accès basé sur les rôles et une connexion unifiée donnent à chaque rôle interne, à chaque partenaire externe et à chaque candidat public exactement la vue et les droits que ses responsabilités exigent. Une plateforme de gestion des ressources humaines est développée comme projet jumeau, et IntelliConnect est conçue pour s'y connecter afin que les données de talents et de recrutement alimentent la même vision partenariale.

Au-delà des opérations quotidiennes, IntelliConnect transforme les données partenariales en aide à la décision. Elle associe des techniques avancées de scoring et de classement à un appariement sémantique vectoriel pour produire des scores de santé clairs, aussi bien pour les partenaires que pour les projets, afin que les analystes puissent repérer les points d'attention avant qu'une situation ne se dégrade. La plateforme applique également les pratiques actuelles d'IA agentique et d'ingénierie de harnais pour offrir un assistant de niveau production qui aide les analystes à produire des rapports exécutifs conformes aux standards de qualité de l'entreprise et exportables en PDF ou en Excel. L'ensemble du système est bâti pour un usage en production, avec un soin constant porté à la fiabilité, à l'observabilité et à la traçabilité du comportement de l'IA, afin que l'intelligence fournie soit digne de confiance dans un contexte d'entreprise réel.

\bigskip
\noindent\textbf{Mots-clés :} gestion des partenariats d'entreprise, gestion de projets, contrôle d'accès basé sur les rôles, scoring de santé des partenaires et des projets, appariement sémantique, IA agentique, génération de rapports, observabilité en production.
\selectlanguage{english}
